%% kapa: exam=1 round=33 year=2022 period=2 form=A subject=accounting subject_ko=회계학 seq=29 number=69
%% source: Q-Net 감정평가사 기출문제 (cst00309) artlSeq=5209997 file=raw/1008/5209997/2022년도 제33회 감정평가사 1차시험 2교시(감정평가 관계법규, 회계학).pdf
%% conversion: kordoc 4.24.51 -> tools/gen1.py (auto); review=proofread-sonnet; flags=none
\begin{question}{69}
(주)감평은 확정급여제도를 운영하고 있으며, 20×1년도 관련 자료는 다음과 같다. 20×1년도 기타포괄손익으로 인식할 확정급여채무의 재측정요소는?
\begin{framedtext}
{\small\begin{tabular}{@{}lr@{}}
기초 확정급여채무의 현재가치 & ￦100,000 \\
기초 사외적립자산의 공정가치 & 90,000 \\
퇴직금 지급액(사외적립자산에서 지급) & 12,000 \\
포괄손익계산서 상 당기손익 인식 퇴직급여 관련 비용 & 28,000 \\
이자비용 & 10,000 \\
이자수익 & 9,000 \\
기말 확정급여채무의 현재가치 & 128,000 \\
기말 사외적립자산의 공정가치 & 99,000 \\
\end{tabular}\par}
\end{framedtext}
\begin{choices}
  \item 재측정손실 ￦2,000
  \item 재측정손실 ￦3,000
  \item 재측정손익 없음
  \item 재측정이익 ￦2,000
  \item 재측정이익 ￦3,000
\end{choices}
\end{question}
\answer{2}
